\subsection{Why a triple must be selected using all four states}
\label{sec:triple-caution}
The preceding proof selects its triple according to the middle masses.
The following constant-curvature example explains why that step cannot
be discarded. Write $P_{UV}$ for a two-point Jensen distortion and
$T_{UVW}$ for the three-point distortion, always with the original weights.

\begin{proposition}[The two three-point inequalities alone give no bound]
Let $q=P_{BC}$, $S=P_{AB}+P_{BC}$, and $T=T_{ABC}$. There is no finite
degree-only bound on a parameter $R$ from the inequalities
\[
 q\leq T/R^2,\qquad S\leq2T/R
\]
for arbitrary positive triple weights, even if the curvature is the
constant polynomial $g\equiv1$ and the locations are fixed.
\end{proposition}
\begin{proof}
Take $\varphi(x)=x^2/2$, $(A,B,C)=(0,1/2,1)$, and original weights
$(1,\varepsilon^2,\varepsilon)$ for $0<\varepsilon\leq1$.
Common normalization of all weights is immaterial. The exact weighted
variance identities give
\[
 P_{AB}=\frac{\varepsilon^2}{8(1+\varepsilon^2)},\qquad
 q=\frac{\varepsilon^2}{8(1+\varepsilon)},\qquad
 T=\frac{\varepsilon+(\varepsilon^2+\varepsilon^3)/4}
          {2(1+\varepsilon+\varepsilon^2)}.
\]
In particular $q\leq\varepsilon^2/8$, $S\leq\varepsilon^2/4$, and
$T\geq\varepsilon/6$. Put $R=\varepsilon^{-1/2}$. Then
\[
 q\leq\varepsilon^2/8\leq\varepsilon T=T/R^2,
 \qquad
 S\leq\varepsilon^2/4\leq\varepsilon^{3/2}/3
    \leq2\sqrt\varepsilon\,T=2T/R.
\]
As $\varepsilon\downarrow0$, $R\to\infty$ while the degree remains
zero and all locations remain fixed.
\end{proof}


Here $R$ is only a formal parameter, not the price of a four-state
hierarchy. The example obstructs the isolated three-point inference.
The proof of Theorem~\ref{thm:weighted-sharp} instead uses both triples
through the common denominator $Q$ and selects one with $u\le v$.
